\section{Integración y continuación numérica}
\label{sec:integracion_continuacion}

\HAFOCardContinuation

El procedimiento numérico distingue tres operaciones:
integrar un problema de valor inicial de Caputo, continuar el sistema no suave
sin borrar su historia y reiniciar después el sistema autónomo para evaluar una
condición inicial bajo un convenio común de memoria.

\subsection{Esquema ABM--PECE de memoria completa}
\label{sec:abm}

Sea
\begin{equation}
 {}^{\mathrm C}D_{t_0}^{q}X(t)=F(t,X(t)),\qquad X(t_0)=X_0,
 \qquad 0<q\leq 1,
\label{eq:pvi_abm}
\end{equation}
y sea \(t_n=t_0+nh\) una malla uniforme. La forma de Volterra
\cref{eq:volterra} se evalúa en \(t_{n+1}\):
\begin{equation}
 X(t_{n+1})=X_0+
 \frac{1}{\Gamma(q)}
 \sum_{j=0}^{n}
 \int_{t_j}^{t_{j+1}}
 (t_{n+1}-\tau)^{q-1}F(\tau,X(\tau))\dd\tau.
\label{eq:volterra_por_intervalos}
\end{equation}
El esquema predictor--corrector de Adams--Bashforth--Moulton (ABM--PECE)
se obtiene al aproximar \(F\) por una constante en el predictor y por una
interpolación lineal en el corrector \cite{DiethelmFordFreed2002}.

En el predictor se sustituye
\(F(\tau,X(\tau))\) por \(F_j=F(t_j,X_j)\) en cada subintervalo. Como
\begin{align}
\int_{t_j}^{t_{j+1}}(t_{n+1}-\tau)^{q-1}\dd\tau
&=\frac{h^q}{q}
\left[(n+1-j)^q-(n-j)^q\right],
\end{align}
y \(q\Gamma(q)=\Gamma(q+1)\), resulta
\begin{equation}
 X_{n+1}^{\mathrm P}=X_0+
 \frac{h^q}{\Gamma(q+1)}
 \sum_{j=0}^{n}b_{j,n+1}F_j,
\qquad
 b_{j,n+1}=(n+1-j)^q-(n-j)^q.
\label{eq:abm_predictor}
\end{equation}

Para el corrector se interpola en cada intervalo mediante
\[
 F(t_j+hu,X(t_j+hu))\approx(1-u)F_j+uF_{j+1},\qquad 0\leq u\leq1.
\]
Al poner \(r=n+1-j\), su contribución a la integral, antes de dividir por
\(\Gamma(q)\), es
\(h^q[L_q(r)F_j+R_q(r)F_{j+1}]\), donde
\begin{align}
 I_q(r)&=\int_0^1(r-u)^{q-1}\,\dd u
       =\frac{r^q-(r-1)^q}{q},\\
 R_q(r)&=\int_0^1u(r-u)^{q-1}\,\dd u
       =rI_q(r)-\frac{r^{q+1}-(r-1)^{q+1}}{q+1},\\
 L_q(r)&=I_q(r)-R_q(r).
 \label{eq:abm_integrales_pesos}
\end{align}
La segunda igualdad se obtiene con \(v=r-u\), de modo que
\(u=r-v\). Al llevar las expresiones al denominador \(q(q+1)\),
\begin{align}
 q(q+1)L_q(r)&=(q+1-r)r^q+(r-1)^{q+1},\\
 q(q+1)R_q(r)&=r^{q+1}-(r+q)(r-1)^q.
\end{align}
El nodo inicial recibe sólo \(L_q(n+1)\); cada nodo interior \(j\) recibe
\(R_q(n+2-j)+L_q(n+1-j)\); el extremo nuevo recibe \(R_q(1)\).
Por tanto,
\begin{align}
 q(q+1)L_q(n+1)&=n^{q+1}-(n-q)(n+1)^q,\\
 q(q+1)[R_q(r+1)+L_q(r)]
  &=(r+1)^{q+1}-2r^{q+1}+(r-1)^{q+1},\\
 q(q+1)R_q(1)&=1.
\end{align}
Finalmente, \(q(q+1)\Gamma(q)=\Gamma(q+2)\). El valor desconocido
\(F_{n+1}\) del extremo se evalúa en el predictor y se obtiene
\begin{equation}
\begin{aligned}
 X_{n+1}
 =X_0+\frac{h^q}{\Gamma(q+2)}
 \left[
 F\!\left(t_{n+1},X_{n+1}^{\mathrm P}\right)
 +\sum_{j=0}^{n}a_{j,n+1}F_j
 \right],
\end{aligned}
\label{eq:abm_corrector}
\end{equation}
con
\begin{equation}
a_{j,n+1}=
\begin{cases}
n^{q+1}-(n-q)(n+1)^q, & j=0,\\[0.25em]
(n-j+2)^{q+1}+(n-j)^{q+1}
-2(n-j+1)^{q+1}, & 1\leq j\leq n.
\end{cases}
\label{eq:pesos_abm}
\end{equation}
La secuencia \emph{PECE} significa: predecir \(X_{n+1}^{\mathrm P}\), evaluar
el campo en ese predictor, corregir para obtener \(X_{n+1}\) y evaluar
finalmente \(F_{n+1}\), que se conserva para el siguiente paso.

Las sumas de \cref{eq:abm_predictor,eq:abm_corrector} comienzan siempre en
\(j=0\). Por ello, una integración de \(N\) pasos utiliza toda la historia
\(\{F_0,\ldots,F_N\}\). El término \emph{memoria completa} indica que ningún
término anterior se elimina ni se sustituye por una
ventana móvil.

\begin{remark}[Recuperación del caso entero]
Para \(q=1\), \(b_{j,n+1}=1\). Además, los pesos del corrector son \(1\) en
los extremos y \(2\) en los nodos interiores; el factor
\(\Gamma(3)^{-1}=1/2\) produce la regla trapezoidal compuesta. Así,
\cref{eq:abm_predictor,eq:abm_corrector} recuperan las cuadraturas compuestas
rectangular izquierda y trapezoidal de la ecuación integral; el extremo
derecho del corrector se evalúa en el predictor. En particular,
\[
 X_{n+1}^{\mathrm P}=X_0+h\sum_{j=0}^{n}F_j
\]
satisface \(X_{n+1}^{\mathrm P}=X_n^{\mathrm P}+hF_n\) para \(n\geq1\).
El corrector, en cambio, satisface
\[
 X_n=X_0+\frac h2\left[F_0+2\sum_{j=1}^{n-1}F_j
                    +F(t_n,X_n^{\mathrm P})\right],
 \qquad
 X_n^{\mathrm P}-X_n=\frac h2[F_0-F(t_n,X_n^{\mathrm P})].
\]
En consecuencia, el predictor no coincide, en general, con el paso local de
Adams--Bashforth de orden uno
\(X_n+hF_n\) iniciado en el estado corregido. Por ejemplo, para
\(\dot x=x\), \(x_0=1\), se obtiene \(x_1=1+h+h^2/2\) y
\(x_2^{\mathrm P}-(x_1+hx_1)=-h^2/2\). La identificación con cuadraturas
integrales conserva los pesos del esquema y precisa su límite entero.
\end{remark}

\subsection{Regularidad del arranque y consistencia de la cuadratura}
\label{subsec:numerical-startup-consistency}

La integración exacta del núcleo singular no elimina la singularidad temporal
de la solución. Supóngase que \(F\) es \(C^1\) cerca de \((t_0,X_0)\) y que
la solución permanece inicialmente en ese entorno. La acotación del campo en
\cref{eq:volterra} implica
\[
 \norm{X(t_0+s)-X_0}\leq \frac{M s^q}{\Gamma(q+1)}.
\]
Como \(0<q\leq1\), el incremento temporal \(s\) es también \(O(s^q)\) cuando
\(s\to0^+\). La regularidad de \(F\) da entonces
\(F(t_0+s,X(t_0+s))=F_0+R(s)\), con \(F_0=F(t_0,X_0)\) y
\(\norm{R(s)}\leq C s^q\). Al sustituir en Volterra, el término constante
se integra exactamente. Para el resto, el cambio \(u=s v\) produce
\[
 \begin{aligned}
 \frac1{\Gamma(q)}\int_0^s(s-u)^{q-1}\norm{R(u)}\,\dd u
 &\leq\frac{C s^{2q}}{\Gamma(q)}
        \int_0^1(1-v)^{q-1}v^q\,\dd v\\
 &=\frac{C\Gamma(q+1)}{\Gamma(2q+1)}s^{2q}.
 \end{aligned}
\]
Por tanto,
\begin{equation}
 X(t_0+s)=X_0+\frac{s^q}{\Gamma(q+1)}F_0+O(s^{2q}).
 \label{eq:numerical-startup-expansion}
\end{equation}
Si \(0<q<1\) y \(F_0\ne0\), dividir esta identidad por \(s\) muestra que
\(\norm{X(t_0+s)-X_0}/s\) no es acotado al aproximarse al terminal inferior:
el término principal es de orden \(s^{q-1}\) y el resto es menor por un
factor \(s^q\). La suavidad del campo no basta, pues, para suponer derivadas
temporales acotadas de la solución.

Para aislar el error de cuadratura, supóngase provisionalmente que
\(g(t)=F(t,X(t))\) pertenece a \(C^2([t_0,t_0+T])\), con
\(\norm{g'}\leq M_1\) y \(\norm{g''}\leq M_2\). En
\([a,b]=[t_j,t_{j+1}]\), \(b-a=h\), la interpolación constante satisface
\(\norm{g(t)-g(a)}\leq M_1h\). Para la interpolación lineal \(\ell_j\), el
resto \(r=g-\ell_j\) cumple \(r(a)=r(b)=0\) y \(r''=g''\). Integrar dos
veces e imponer ambas condiciones de borde da
\[
 \begin{aligned}
 r(t)={}&-\frac{b-t}{h}\int_a^t(u-a)g''(u)\,\dd u\\
        &-\frac{t-a}{h}\int_t^b(b-u)g''(u)\,\dd u.
 \end{aligned}
\]
La desigualdad triangular, válida también para campos vectoriales, implica
\[
 \begin{aligned}
 \norm{r(t)}
 &\leq\frac{M_2}{2h}\bigl[(b-t)(t-a)^2+(t-a)(b-t)^2\bigr]\\
 &=\frac{M_2}{2}(t-a)(b-t)\leq\frac{M_2h^2}{8}.
 \end{aligned}
\]
Al integrar los errores contra el núcleo positivo y usar
\(\int_0^T u^{q-1}\,\dd u/\Gamma(q)=T^q/\Gamma(q+1)\), se obtienen los
defectos sobre datos nodales exactos
\begin{equation}
 \varepsilon_{\mathrm P}\leq\frac{M_1T^q}{\Gamma(q+1)}h,
 \qquad
 \varepsilon_{\mathrm L}\leq\frac{M_2T^q}{8\Gamma(q+1)}h^2.
 \label{eq:numerical-quadrature-defects}
\end{equation}
El primero corresponde al predictor y el segundo a la interpolación lineal
con el valor exacto del campo en el extremo nuevo. Si \(F\) es
\(L\)-Lipschitz respecto del estado en un entorno que contiene la solución y
el predictor, evaluar ese extremo en \(X_{n+1}^{\mathrm P}\) añade a lo sumo
\(L\varepsilon_{\mathrm P}\) al error del campo. Su peso en el corrector es
\(h^q/\Gamma(q+2)\); por ello el defecto PECE satisface
\begin{equation}
 \varepsilon_{\mathrm{PECE}}
 \leq \varepsilon_{\mathrm L}
       +\frac{Lh^q}{\Gamma(q+2)}\varepsilon_{\mathrm P}
 =O(h^2+h^{1+q}).
 \label{eq:numerical-pece-consistency}
\end{equation}
Esta es una estimación de consistencia al insertar valores exactos en la
cuadratura. Un orden global requiere además controlar la propagación de los
errores nodales \cite{DiethelmFordFreed2002}. Cuando \(g'\) o \(g''\) no son
acotadas cerca de \(t_0\), no existen las constantes uniformes usadas en
\cref{eq:numerical-quadrature-defects}. En ese caso se debe estimar por
separado el tramo inicial, introducir correcciones de arranque o analizar
una malla graduada. Los refinamientos de \cref{sec:control_numerico} se
interpretan según la regularidad efectiva de cada problema, sin atribuirles
automáticamente el orden del régimen suave.

\subsection{Continuación causal del Chua no suave}
\label{sec:continuacion_causal}

La semilla sesgada obtenida mediante
\cref{eq:cierre_dc_sesgado,eq:semilla_sesgada} pertenece a un sistema
auxiliar linealizado. Antes de discretizar el parámetro de continuación se
deriva la familia afín y se comprueba que conserva la componente continua, el
modo armónico y los dos extremos:
\begingroup
\def\HAFOderivationpart{10}
% Este archivo es una biblioteca de derivaciones. Cada bloque se inserta en el
% lugar del cuerpo donde se usa, mediante \HAFOderivationpart. No se imprime
% como apéndice independiente.
\ifnum\HAFOderivationpart=1\relax
\subsection{Derivación completa de la forma de Lur'e}
\label{app:verificacion_lure}

Sea \(h(x)=\ell x+\psi(x)\). La primera ecuación de los dos modelos es
\[
{}^{\mathrm C}D_{t_0}^{q}x
=\alpha(y-x-h(x)).
\]
Al distribuir \(\alpha\),
\begin{align}
\alpha(y-x-h(x))
&=\alpha y-\alpha x-\alpha\ell x-\alpha\psi(x)\nonumber\\
&=-\alpha(1+\ell)x+\alpha y-\alpha\psi(x).
\label{eq:expansion_primera_fila}
\end{align}
Por otra parte,
\[
PX=
\begin{pmatrix}
-\alpha(1+\ell)x+\alpha y\\
x-y+z\\
-\beta y-\gamma z
\end{pmatrix},
\qquad
b\psi(r^{\T}X)=
\begin{pmatrix}
-\alpha\psi(x)\\0\\0
\end{pmatrix}.
\]
La suma reproduce \cref{eq:expansion_primera_fila} y las otras dos
ecuaciones. Para la saturación se toma
\(\ell=m_1\) y
\(\psi=(m_0-m_1)\operatorname{sat}\); para la arctangente,
\(\ell=a_1\) y \(\psi=a_2\arctan(\rho\,\cdot)\).

\fi
\ifnum\HAFOderivationpart=2\relax
\subsection{Derivación completa del determinante y la transferencia}
\label{app:transferencia_explicita}

Con \(\lambda=s^q\),
\begin{equation}
\lambda I-P=
\begin{pmatrix}
\lambda+\alpha(1+\ell)&-\alpha&0\\
-1&\lambda+1&-1\\
0&\beta&\lambda+\gamma
\end{pmatrix}.
\label{eq:matriz_resolvente}
\end{equation}
El menor de la posición \((1,1)\) es
\begin{align}
Q(\lambda)
&=
\det
\begin{pmatrix}
\lambda+1&-1\\
\beta&\lambda+\gamma
\end{pmatrix}\nonumber\\
&=(\lambda+1)(\lambda+\gamma)+\beta.
\label{eq:menor_Q}
\end{align}
Al desarrollar \(\det(\lambda I-P)\) por la primera fila,
\begin{align}
\Delta_\ell(\lambda)
&=\bigl[\lambda+\alpha(1+\ell)\bigr]Q(\lambda)
-(-\alpha)
\det
\begin{pmatrix}
-1&-1\\
0&\lambda+\gamma
\end{pmatrix}\nonumber\\
&=\bigl[\lambda+\alpha(1+\ell)\bigr]Q(\lambda)
-\alpha(\lambda+\gamma).
\label{eq:determinante_transferencia}
\end{align}

Para calcular la primera componente de
\((\lambda I-P)^{-1}b\), la regla de Cramer sustituye la primera columna de
\cref{eq:matriz_resolvente} por
\(b=(-\alpha,0,0)^{\T}\). El determinante del numerador es
\[
\det
\begin{pmatrix}
-\alpha&-\alpha&0\\
0&\lambda+1&-1\\
0&\beta&\lambda+\gamma
\end{pmatrix}
=-\alpha Q(\lambda).
\]
Como \(r^{\T}\) selecciona esa primera componente,
\begin{equation}
W_q(s)
=r^{\T}(s^qI-P)^{-1}b
=-\alpha\frac{Q(s^q)}{\Delta_\ell(s^q)}.
\label{eq:transferencia_cramer}
\end{equation}

La identidad
\[
P-s^qI=-(s^qI-P)
\]
implica
\[
r^{\T}(P-s^qI)^{-1}b=-W_q(s).
\]
Por ello, cambiar el orden de los términos dentro de la resolvente cambia
simultáneamente el signo de la transferencia y el signo del cierre. No cambia
la solución física si ambas modificaciones se realizan de manera coherente.

\fi
\ifnum\HAFOderivationpart=3\relax
\subsection{Derivación completa de los equilibrios}
\label{app:equilibrios}

En un equilibrio \(e=(x_e,y_e,z_e)^{\T}\), las dos últimas ecuaciones de
ambos sistemas dan
\[
0=x_e-y_e+z_e,\qquad
0=-\beta y_e-\gamma z_e.
\]
De la primera,
\[
z_e=y_e-x_e.
\]
Al sustituirla en la segunda,
\begin{align*}
0
&=-\beta y_e-\gamma(y_e-x_e)\\
&=-(\beta+\gamma)y_e+\gamma x_e.
\end{align*}
Si \(\beta+\gamma\neq0\),
\begin{equation}
y_e=\frac{\gamma}{\beta+\gamma}x_e,\qquad
z_e=-\frac{\beta}{\beta+\gamma}x_e.
\label{eq:equilibrios_yz}
\end{equation}
La primera ecuación de equilibrio exige
\[
0=y_e-x_e-h(x_e).
\]
Mediante \cref{eq:equilibrios_yz},
\begin{equation}
h(x_e)+\frac{\beta}{\beta+\gamma}x_e=0.
\label{eq:equilibrio_escalar_comun}
\end{equation}
Esta única ecuación escalar determina todos los equilibrios; las otras dos
coordenadas se recuperan después con \cref{eq:equilibrios_yz}.

\subsubsection{Equilibrios del Chua no suave}

En la región \(|x_e|<1\), \(h(x_e)=m_0x_e\). Entonces
\begin{equation}
\left(m_0+\frac{\beta}{\beta+\gamma}\right)x_e=0.
\label{eq:equilibrio_interior_ns}
\end{equation}
Cuando el coeficiente no es cero, la única raíz interior es \(x_e=0\).

Para \(x_e>1\),
\[
h(x_e)=m_1x_e+(m_0-m_1),
\]
y \cref{eq:equilibrio_escalar_comun} queda
\[
\left(m_1+\frac{\beta}{\beta+\gamma}\right)x_e+(m_0-m_1)=0.
\]
Por tanto,
\begin{equation}
x_\star=
\frac{m_1-m_0}
{m_1+\dfrac{\beta}{\beta+\gamma}}.
\label{eq:raiz_exterior_ns}
\end{equation}
Si \(x_\star>1\), existe el equilibrio positivo. Como la no linealidad es
impar, existe también la raíz \(-x_\star<-1\). Los tres equilibrios son
\begin{equation}
E_0=(0,0,0),\qquad
E_\pm=
\left(
\pm x_\star,\,
\pm\frac{\gamma}{\beta+\gamma}x_\star,\,
\mp\frac{\beta}{\beta+\gamma}x_\star
\right).
\label{eq:equilibrios_ns_simbolicos}
\end{equation}
La condición regional \(x_\star>1\) debe comprobarse; no basta con resolver la
ecuación lineal exterior.

\subsubsection{Equilibrios del Chua con arctangente}

Al usar
\(h(x)=a_1x+a_2\arctan(\rho x)\) en
\cref{eq:equilibrio_escalar_comun}, se obtiene
\begin{equation}
\left(a_1+\frac{\beta}{\beta+\gamma}\right)x_e
+a_2\arctan(\rho x_e)=0.
\label{eq:equilibrio_escalar_arctan}
\end{equation}
La expresión es impar, de modo que \(x_e=0\) siempre es raíz y toda raíz
positiva no nula tiene una compañera negativa. Las raíces no nulas se
determinan numéricamente y se sustituyen en \cref{eq:equilibrios_yz}. Esta
ecuación se evalúa con \(\rho=1\) para el caso bibliográfico y con el valor de
\(\rho\) de \cref{eq:parametros_c590} para el caso \texttt{c590}; no deben
intercambiarse las dos parametrizaciones.

\fi
\ifnum\HAFOderivationpart=4\relax
\subsection{Derivación completa de los jacobianos regionales}
\label{app:jacobianos}

Para la forma de Lur'e,
\[
F(X)=PX+b\psi(r^{\T}X).
\]
La regla de la cadena produce
\begin{equation}
J_e=P+b\,\psi'(r^{\T}e)r^{\T}.
\label{eq:jacobiano_lure}
\end{equation}
Como \(r^{\T}e=x_e\), solamente cambia la entrada \((1,1)\).

En el sistema no suave,
\[
\psi'_{\mathrm{ns}}(x)=
\begin{cases}
m_0-m_1,&|x|<1,\\
0,&|x|>1.
\end{cases}
\]
Por ello,
\begin{equation}
J_{\mathrm{in}}=
\begin{pmatrix}
-\alpha(1+m_0)&\alpha&0\\
1&-1&1\\
0&-\beta&-\gamma
\end{pmatrix},
\qquad
J_{\mathrm{out}}=
\begin{pmatrix}
-\alpha(1+m_1)&\alpha&0\\
1&-1&1\\
0&-\beta&-\gamma
\end{pmatrix}.
\label{eq:jacobianos_ns}
\end{equation}
En \(x=\pm1\) la derivada clásica no existe. Los equilibrios usados en este
reporte no se encuentran sobre esas superficies, por lo que se emplea el
jacobiano de la región que contiene a cada uno.

Para la arctangente,
\[
\psi'_{\mathrm a}(x)
=\frac{a_2\rho}{1+\rho^2x^2},
\]
y
\begin{equation}
J_e=
\begin{pmatrix}
\displaystyle
-\alpha\left(1+a_1+
\frac{a_2\rho}{1+\rho^2x_e^2}\right)
&\alpha&0\\[3mm]
1&-1&1\\
0&-\beta&-\gamma
\end{pmatrix}.
\label{eq:jacobiano_arctan}
\end{equation}
Los valores propios de
\cref{eq:jacobianos_ns,eq:jacobiano_arctan} se clasifican con
\cref{eq:matignon}; la estabilidad local no sustituye el sondeo de cuenca.

\fi
\ifnum\HAFOderivationpart=5\relax
\subsection{Derivación completa del polinomio característico}
\label{app:polinomio_caracteristico}

Sea \(g=h'(x_e)\) la pendiente total de la no linealidad en un equilibrio.
El jacobiano de cualquiera de las dos familias puede escribirse como
\begin{equation}
J(g)=
\begin{pmatrix}
-\alpha(1+g)&\alpha&0\\
1&-1&1\\
0&-\beta&-\gamma
\end{pmatrix}.
\label{eq:jacobiano_pendiente_total}
\end{equation}
Por tanto,
\[
\lambda I-J(g)=
\begin{pmatrix}
\lambda+\alpha(1+g)&-\alpha&0\\
-1&\lambda+1&-1\\
0&\beta&\lambda+\gamma
\end{pmatrix}.
\]
El desarrollo por la primera fila, igual al de
\cref{eq:determinante_transferencia}, produce
\begin{equation}
\chi_g(\lambda)
=\bigl[\lambda+\alpha(1+g)\bigr]
\bigl[(\lambda+1)(\lambda+\gamma)+\beta\bigr]
-\alpha(\lambda+\gamma).
\label{eq:polinomio_caracteristico_agrupado}
\end{equation}
Como
\[
(\lambda+1)(\lambda+\gamma)+\beta
=\lambda^2+(1+\gamma)\lambda+(\beta+\gamma),
\]
la multiplicación y agrupación por potencias de \(\lambda\) da
\begin{equation}
\begin{aligned}
\chi_g(\lambda)
={}&\lambda^3+
\bigl[1+\gamma+\alpha(1+g)\bigr]\lambda^2\\
&+\bigl[
\beta+\gamma+\alpha(1+g)(1+\gamma)-\alpha
\bigr]\lambda\\
&+\alpha\bigl[(1+g)(\beta+\gamma)-\gamma\bigr].
\end{aligned}
\label{eq:polinomio_caracteristico_expandido}
\end{equation}
En el sistema no suave se sustituye \(g=m_0\) en \(E_0\) y \(g=m_1\)
en \(E_\pm\). En el sistema con arctangente,
\[
g=a_1+\frac{a_2\rho}{1+\rho^2x_e^2}.
\]
Las raíces de \cref{eq:polinomio_caracteristico_expandido} son los valores
propios utilizados en \cref{eq:matignon}; así, la clasificación local se
obtiene a partir del mismo campo que define la integración.

\fi
\ifnum\HAFOderivationpart=6\relax
\subsection{Derivación completa para la saturación centrada}
\label{app:df_saturacion_centrada}

Sea
\[
\psi_{\mathrm{ns}}(\sigma)=\Delta\operatorname{sat}(\sigma),
\qquad \Delta=m_0-m_1.
\]
Si \(0<A\leq1\), no hay saturación y
\[
\psi_{\mathrm{ns}}(A\sin\theta)=\Delta A\sin\theta.
\]
Como \(\int_0^\pi\sin^2\theta\dd\theta=\pi/2\),
\[
N_{\mathrm{ns}}(A)
=\frac{2}{\pi A}\Delta A\frac{\pi}{2}
=\Delta.
\]

Para \(A>1\), defínase
\[
\theta_0=\arcsin\left(\frac1A\right),
\qquad 0<\theta_0<\frac{\pi}{2}.
\]
En \(0\leq\theta\leq\theta_0\) y
\(\pi-\theta_0\leq\theta\leq\pi\), la entrada permanece en la parte lineal;
en el intervalo central la saturación vale \(1\). Entonces
\begin{align}
N_{\mathrm{ns}}(A)
&=\frac{2\Delta}{\pi A}
\left[
2A\int_0^{\theta_0}\sin^2\theta\dd\theta
+\int_{\theta_0}^{\pi-\theta_0}\sin\theta\dd\theta
\right]\nonumber\\
&=\frac{2\Delta}{\pi A}
\left[
A\theta_0-\frac{A}{2}\sin(2\theta_0)
+2\cos\theta_0
\right]\nonumber\\
&=\frac{2\Delta}{\pi A}
\left[A\theta_0+\cos\theta_0\right].
\label{eq:df_sat_pasos}
\end{align}
En el último paso se usó
\[
\frac{A}{2}\sin(2\theta_0)
=A\sin\theta_0\cos\theta_0
=\cos\theta_0.
\]
Finalmente,
\[
\cos\theta_0
=\sqrt{1-\frac1{A^2}}
=\frac{\sqrt{A^2-1}}{A},
\]
y \cref{eq:df_sat_pasos} se convierte en
\[
N_{\mathrm{ns}}(A)
=\Delta\frac{2}{\pi}
\left[
\arcsin\left(\frac1A\right)
+\frac{\sqrt{A^2-1}}{A^2}
\right],
\]
que coincide con \cref{eq:df_saturacion}.

\fi
\ifnum\HAFOderivationpart=7\relax
\subsection{Derivación completa de la saturación sesgada}
\label{app:df_saturacion_sesgada}

Considérese
\[
u(\theta)=c+A\cos\theta,\qquad A>0.
\]
Para escribir una fórmula que incluya los casos parcialmente y completamente
saturados, se define
\[
[v]_{-1}^{1}=\max(-1,\min(1,v))
\]
y los ángulos
\begin{equation}
\theta_+
=\arccos\left[\frac{1-c}{A}\right]_{-1}^{1},
\qquad
\theta_-
=\arccos\left[\frac{-1-c}{A}\right]_{-1}^{1}.
\label{eq:angulos_saturacion_sesgada}
\end{equation}
Como \((1-c)/A>(-1-c)/A\) y \(\arccos\) es decreciente,
\(0\leq\theta_+\leq\theta_-\leq\pi\). En el intervalo
\([0,\pi]\),
\[
\operatorname{sat}(u(\theta))=
\begin{cases}
1,&0\leq\theta\leq\theta_+,\\
c+A\cos\theta,&\theta_+\leq\theta\leq\theta_-,\\
-1,&\theta_-\leq\theta\leq\pi.
\end{cases}
\]
La simetría respecto de \(2\pi\) permite integrar solamente en
\([0,\pi]\). Para \(\psi=\Delta\operatorname{sat}\),
\begin{align}
\Psi_0(c,A)
=\frac{\Delta}{\pi}\Big[
&\theta_+
+c(\theta_--\theta_+)
+A(\sin\theta_--\sin\theta_+)\nonumber\\
&-(\pi-\theta_-)
\Big].
\label{eq:psi0_saturacion_cerrada}
\end{align}

Para la componente fundamental se usan
\[
\int\cos\theta\dd\theta=\sin\theta,\qquad
\int\cos^2\theta\dd\theta
=\frac{\theta}{2}+\frac{\sin2\theta}{4}.
\]
La sustitución por intervalos da
\begin{align}
\Psi_1(c,A)
=\frac{2\Delta}{\pi}\Bigg[
&\sin\theta_++\sin\theta_-
+c(\sin\theta_--\sin\theta_+)\nonumber\\
&+\frac{A}{2}(\theta_--\theta_+)
+\frac{A}{4}
\bigl(\sin2\theta_--\sin2\theta_+\bigr)
\Bigg].
\label{eq:psi1_saturacion_cerrada}
\end{align}
Las operaciones de recorte en
\cref{eq:angulos_saturacion_sesgada} hacen que
\cref{eq:psi0_saturacion_cerrada,eq:psi1_saturacion_cerrada} sigan siendo
válidas cuando alguna de las tres regiones tiene longitud cero.

La ecuación de media puede reducirse a una igualdad escalar. De
\cref{eq:transferencia_cramer},
\begin{align}
W_q(0)
&=-\alpha
\frac{\beta+\gamma}
{\alpha(1+\ell)(\beta+\gamma)-\alpha\gamma}\nonumber\\
&=-\frac{\beta+\gamma}
{(1+\ell)(\beta+\gamma)-\gamma}.
\label{eq:ganancia_dc}
\end{align}
Por tanto,
\begin{equation}
c=W_q(0)\Psi_0(c,A),
\label{eq:cierre_dc_escalar}
\end{equation}
y el cierre armónico es
\[
1-W_q(\ii\omega)\frac{\Psi_1(c,A)}{A}=0.
\]
Estas tres ecuaciones reales ---una para la media, y las partes real e
imaginaria del cierre--- determinan \((c,A,\omega)\).

\fi
\ifnum\HAFOderivationpart=8\relax
\subsection{Derivación completa para la arctangente}
\label{app:df_arctan}

Sea
\[
\psi_{\mathrm a}(\sigma)=a_2\arctan(\rho\sigma),
\qquad
\xi=\rho A.
\]
Se define
\[
I(\xi)=
\int_0^\pi\arctan(\xi\sin\theta)\sin\theta\dd\theta.
\]
Como \(I(0)=0\), puede calcularse \(I\) integrando su derivada:
\begin{align}
I'(\xi)
&=\int_0^\pi
\frac{\sin^2\theta}{1+\xi^2\sin^2\theta}\dd\theta\nonumber\\
&=\frac1{\xi^2}
\left[
\pi-\int_0^\pi
\frac{\dd\theta}{1+\xi^2\sin^2\theta}
\right]\nonumber\\
&=\frac{\pi}{\xi^2}
\left(1-\frac1{\sqrt{1+\xi^2}}\right).
\label{eq:derivada_integral_arctan}
\end{align}
En la última igualdad se usó la integral estándar
\[
\int_0^\pi\frac{\dd\theta}{1+\xi^2\sin^2\theta}
=\frac{\pi}{\sqrt{1+\xi^2}}.
\]
Una primitiva que se anula en \(\xi=0\) es
\begin{equation}
I(\xi)
=\pi\frac{\sqrt{1+\xi^2}-1}{\xi}.
\label{eq:integral_arctan}
\end{equation}
Al sustituir \(\xi=\rho A\) en
\cref{eq:df_centrada},
\begin{align}
N_{\mathrm a}(A)
&=\frac{2a_2}{\pi A}I(\rho A)\nonumber\\
&=\frac{2a_2}{\rho A^2}
\left(\sqrt{1+\rho^2A^2}-1\right)\nonumber\\
&=\frac{2a_2\rho}
{1+\sqrt{1+\rho^2A^2}}.
\label{eq:df_arctan_derivada}
\end{align}

Si el cierre exige \(N_{\mathrm a}(A)=k\), la amplitud puede despejarse.
De la forma racionalizada,
\[
1+\sqrt{1+\rho^2A^2}=\frac{2a_2\rho}{k}.
\]
Al aislar la raíz, elevar al cuadrado y simplificar,
\begin{equation}
A^2
=\frac{4a_2(a_2\rho-k)}{k^2\rho}.
\label{eq:amplitud_arctan}
\end{equation}
La solución se acepta solamente si el lado derecho es positivo y la igualdad
previa a elevar al cuadrado tiene lado derecho no menor que \(2\). De manera
equivalente, para \(A>0\), \(k\) tiene el mismo signo que \(a_2\) y se
encuentra estrictamente entre \(0\) y \(a_2\rho\), con el orden de los
extremos ajustado al signo de \(a_2\).

\fi
\ifnum\HAFOderivationpart=9\relax
\subsection{Derivación completa de la identidad modal y su normalización}
\label{app:identidad_modal}

Sea \(P_0=P+bkr^{\T}\). Entonces
\[
\zeta I-P_0=(\zeta I-P)-bkr^{\T}.
\]
Si \(\zeta I-P\) es invertible, el lema del determinante matricial
\(\det(M+uv^{\T})=\det(M)(1+v^{\T}M^{-1}u)\), con
\[
M=\zeta I-P,\qquad u=-bk,\qquad v=r,
\]
produce
\begin{align}
\det(\zeta I-P_0)
&=\det(\zeta I-P)
\left[1-kr^{\T}(\zeta I-P)^{-1}b\right]\nonumber\\
&=\det(\zeta I-P)\bigl[1-kG(\zeta)\bigr],
\label{eq:lema_determinante_aplicado}
\end{align}
donde
\[
G(\zeta)=r^{\T}(\zeta I-P)^{-1}b,
\qquad W_q(s)=G(s^q).
\]
Por continuidad, la identidad polinómica se extiende también a los valores en
los que la resolvente de \(P\) es singular. Si
\[
1-kG(\zeta_0)=0,
\]
entonces \(\det(\zeta_0I-P_0)=0\), de modo que existe
\(v\neq0\) con \(P_0v=\zeta_0v\). Siempre que \(r^{\T}v\neq0\), se reemplaza
\(v\) por \(v/(r^{\T}v)\) y se obtiene \(r^{\T}v=1\). Ésta es la
normalización usada en \cref{eq:modo_normalizado,eq:semilla_modal}.

En la rama sesgada,
\[
X_1=\bigl[(\ii\omega_0)^qI-P\bigr]^{-1}b\Psi_1.
\]
Al premultiplicar por \(r^{\T}\),
\[
r^{\T}X_1
=W_q(\ii\omega_0)\Psi_1
=W_q(\ii\omega_0)N_1A.
\]
El cierre \(1-W_qN_1=0\) prueba algebraicamente que
\(r^{\T}X_1=A\).

\fi
\ifnum\HAFOderivationpart=10\relax
\subsection{Derivación completa de la homotopía afín}
\label{app:homotopia_afin}

La aproximación sesgada separa
\[
\psi(\sigma)
=\Psi_0+N_1(\sigma-c)
+\bigl[\psi(\sigma)-\Psi_0-N_1(\sigma-c)\bigr].
\]
Se definen
\begin{equation}
P_{\mathrm a}=P+N_1br^{\T},\qquad
d_{\mathrm a}=b(\Psi_0-N_1c).
\label{eq:parametros_homotopia_afin}
\end{equation}
La familia de campos es
\begin{equation}
\begin{aligned}
F_\varepsilon(X)
&=P_{\mathrm a}X+d_{\mathrm a}\\
&\quad+\varepsilon b
\left[
\psi(r^{\T}X)-\Psi_0-N_1(r^{\T}X-c)
\right],
\qquad 0\leq\varepsilon\leq1.
\end{aligned}
\label{eq:homotopia_afin}
\end{equation}

En \(\varepsilon=0\), la media \(\bar X\) de
\cref{eq:cierre_dc_sesgado} es un equilibrio del auxiliar:
\begin{align*}
F_0(\bar X)
&=P\bar X+bN_1c+b\Psi_0-bN_1c\\
&=P\bar X+b\Psi_0=0.
\end{align*}
Para una perturbación \(\xi=X-\bar X\),
\[
F_0(\bar X+\xi)=P_{\mathrm a}\xi,
\]
de modo que el auxiliar conserva tanto la media como el modo armónico.

En \(\varepsilon=1\), al agrupar términos,
\begin{align*}
F_1(X)
&=PX+bN_1r^{\T}X+b\Psi_0-bN_1c\\
&\quad+b\psi(r^{\T}X)-b\Psi_0-bN_1r^{\T}X+bN_1c\\
&=PX+b\psi(r^{\T}X).
\end{align*}
Así, el extremo final de \cref{eq:homotopia_afin} es exactamente el sistema
físico de \cref{eq:lure_comun}. Una homotopía centrada de la forma
\(P_0X+\varepsilon b[\psi(\sigma)-k\sigma]\) no conserva, en general, la
media \(c\neq0\); por eso no se utiliza para la rama sesgada.
\fi

\endgroup

Con la notación \(\eta=\varepsilon\),
\(\eta=0\) es el auxiliar usado para construir la semilla y \(\eta=1\) es
exactamente el Chua no suave de \cref{eq:chua_no_suave}; no queda un término
residual de la aproximación armónica.

Se usan los niveles
\begin{equation}
\eta_k=0.05k,\qquad k=0,\ldots,20.
\label{eq:niveles_eta}
\end{equation}
Cada nivel ocupa \(60\) unidades de tiempo: las primeras \(30\) permiten el
ajuste al nuevo campo y las siguientes \(30\) registran la respuesta. Si
\(t_k=t_0+60k\), el parámetro se define como
\begin{equation}
\eta(t)=\eta_k,
\qquad t\in[t_k,t_{k+1}),\qquad k=0,\ldots,20.
\label{eq:eta_tiempo}
\end{equation}
Todo el intervalo \([t_0,t_{21}]\) se integra como un solo problema
\begin{equation}
{}^{\mathrm C}D_{t_0}^{q}X(t)=F_{\eta(t)}(X(t)).
\label{eq:continuacion_unica_historia}
\end{equation}
Cuando \(k\) aumenta, únicamente cambia el campo \(F_{\eta_k}\); el índice
inicial de las sumas ABM continúa siendo \(j=0\). Esta es la adaptación
fundamental de la continuación de Leonov--Kuznetsov al operador de Caputo.
Usar el último estado como nueva condición inicial en cada nivel produciría
veintiún problemas diferentes y perdería la memoria acumulada.

La diferencia puede verse directamente en la ecuación integral. Para
\(t\in[t_k,t_{k+1})\), defínase la contribución heredada
\[
 H_k(t)=\frac1{\Gamma(q)}\int_{t_0}^{t_k}
 (t-s)^{q-1}F_{\eta(s)}(X(s))\,\dd s.
\]
Como \(X(t_k)=X_0+H_k(t_k)\), resulta
\begin{equation}
 X(t)=X(t_k)+H_k(t)-H_k(t_k)
 +\frac1{\Gamma(q)}\int_{t_k}^{t}(t-s)^{q-1}
                 F_{\eta_k}(X(s))\,\dd s.
 \label{eq:continuacion_memoria_heredada}
\end{equation}
Un reinicio en \(t_k\) elimina \(H_k(t)-H_k(t_k)\). Para \(q=1\), el
núcleo es constante y esa diferencia es cero. Para \(0<q<1\), el núcleo
pondera de nuevo toda la historia al avanzar \(t\); en general la diferencia
permanece y no puede determinarse sólo a partir de \(X(t_k)\).

\paragraph{Interfaces entre niveles de continuación.}
\label{subsec:numerical-continuation-interfaces}
Si \(\eta(t)\) cambia por saltos, el campo dependiente del tiempo es continuo
por tramos. La formulación integral conserva sentido, pero las cotas de
interpolación de \cref{subsec:numerical-startup-consistency} no pueden
atravesar esos saltos con una derivada temporal acotada. Los tiempos \(t_k\)
deben ser fronteras de los subintervalos de cuadratura. La trayectoria es
continua y posee un único valor \(X(t_k)\), mientras que los valores límite
del campo son, en general, distintos:
\[
 g(t_k^-)=F_{\eta_{k-1}}(X(t_k)),\qquad
 g(t_k^+)=F_{\eta_k}(X(t_k)).
\]
La interpolación lineal en el intervalo que termina en \(t_k\) utiliza el
primer valor; la del intervalo que comienza allí utiliza el segundo. El
valor del integrando en un punto aislado no cambia la integral exacta, pero
su elección como dato nodal sí cambia el interpolante en un intervalo de
longitud positiva. Si una discretización emplea un único valor en esa
interfaz, el defecto correspondiente forma parte del control numérico de la
transición. Esta condición de cuadratura se distingue de conservar la
historia: ambas son necesarias para aplicar el análisis de consistencia
por tramos. El efecto de omitir por completo la historia se calcula en
\cref{subsec:numerical-adm-history}.

La búsqueda sesgada entrega
\begin{equation}
\begin{aligned}
A&=4.5778827210,&
c&=2.7763970033,&
\omega&=2.0402860511,\\
X_{\mathrm{sem}}&=
(7.3542797243,\;0.2909687788,\;-9.3160548187)^{\T}.
\end{aligned}
\label{eq:datos_semilla_no_suave}
\end{equation}
Al completar \(\eta=1\), el último vector de la historia es
\begin{equation}
X_{\mathrm{fin}}=
(0.8223107098,\;1.4652727040,\;1.2923082502)^{\T}.
\label{eq:endpoint_no_suave}
\end{equation}

\subsection{Reinicio controlado para la trayectoria objetivo}
\label{sec:reinicio_objetivo}

La continuación de \cref{eq:continuacion_unica_historia} proporciona un camino
causal desde la semilla auxiliar hasta el sistema original. Para
comparar esa condición final con las condiciones iniciales de los sondeos se
realiza después una operación distinta:
\begin{equation}
{}^{\mathrm C}D_{0}^{q}X(t)=PX(t)+b\psi(r^{\T}X(t)),
\qquad X(0)=X_{\mathrm{fin}}.
\label{eq:reset_objetivo_no_suave}
\end{equation}
El reinicio define una trayectoria autónoma con el mismo instante inicial y
el mismo convenio de memoria de los sondeos. Su segmento postransitorio
constituye la nube de referencia \(\target\) del clasificador.

Así, las dos operaciones tienen propósitos diferentes:
\begin{enumerate}
 \item la continuación conserva la historia para trasladar la semilla sin
 cambiar el problema fraccionario durante el recorrido;
 \item el reinicio estandariza el límite inferior y permite comparar cuencas
 bajo un único convenio de Caputo.
\end{enumerate}

\subsection{Refinamiento fraccionario del caso con arctangente}
\label{sec:refinamiento_c590}

El caso \texttt{c590} de \cref{eq:parametros_c590} sigue una ruta distinta.
Primero se explora el sistema entero \(q=1\), donde una trayectoria queda
determinada por el estado instantáneo y puede integrarse eficientemente con
RK4. La exploración acotada produce un sistema y una trayectoria fuente. De
esa trayectoria se extraen estados postransitorios, que se usan como banco de
semillas para problemas fraccionarios.

Para cada orden \(q\) y cada semilla \(X_0^{(j)}\) se resuelve un problema
independiente:
\begin{equation}
{}^{\mathrm C}D_{0}^{q}X^{(j,q)}(t)
=F_{\mathrm a}\!\left(X^{(j,q)}(t)\right),
\qquad X^{(j,q)}(0)=X_0^{(j)}.
\label{eq:pvi_refinamiento_c590}
\end{equation}
Los órdenes ensayados son
\begin{equation}
q\in\{0.9995,\;0.9998,\;0.9999,\;0.99995\}.
\label{eq:ordenes_refinamiento}
\end{equation}
Cambiar \(q\) modifica el operador y el núcleo de memoria; por esa razón no se
transporta una historia entre órdenes. Cada fila de la exploración es un PVI
completo con su propio \(t_0=0\).

Los candidatos que permanecen acotados se comparan en
\(h\in\{0.0025,0.005,0.01\}\). La regla de selección exige acotamiento en los
tres pasos y un estadístico \(0\)-\(1\) mediano mayor que \(0.8\) en
\(h=0.0025\) y \(h=0.005\). La semilla seleccionada, identificada como
\texttt{seed9}, corresponde al tiempo nominal \(t=240\) de la trayectoria
fuente:
\begin{equation}
X_{0,\mathrm{c590}}=
(5.8642449791,\;1.5847111486,\;3.2155806478)^{\T}.
\label{eq:semilla_c590}
\end{equation}
La trayectoria final se calcula con \(q=0.9999\), \(h=0.005\) y memoria
completa. La exploración entera localiza una región prometedora; la integración
de \cref{eq:pvi_refinamiento_c590} verifica la respuesta bajo el operador de
Caputo que define el sistema estudiado.

\subsection{Control de convergencia numérica}
\label{sec:control_numerico}

La aceptación de convergencia respecto del paso requiere conservar el resultado
al cambiar \(h\) dentro del conjunto declarado y completar las integraciones
sin valores no finitos ni desbordamiento. Este requisito se distingue de la
robustez frente a perturbaciones de la condición inicial. Para las trayectorias
objetivo se registran, como mínimo:
\begin{equation}
\left(
q,\ h,\ T,\ T_{\mathrm{burn}},\ X_0,\
\text{integrador},\ \text{política de memoria}
\right).
\label{eq:contrato_numerico_minimo}
\end{equation}
El estado postransitorio se analiza con la misma especificación utilizada para crear
la nube objetivo. En el caso no suave se integraron \(36\) perturbaciones de
la condición inicial a \(h=0.01\), con ABM y memoria completa: \(12\) para
cada radio \(10^{-6},10^{-4},10^{-2}\). Las \(36\) hicieron contacto con la
nube objetivo. Este resultado documenta robustez respecto de \(X_0\) a ese
paso. La convergencia respecto de \(h\) requiere todavía una comparación
\(h,h/2,h/4\) del mismo problema de valor inicial.

La comparación conservará el campo, los
parámetros, \(X_0\), terminal inferior, horizonte, política de memoria y
observables. Se evaluará el error de trayectoria en una ventana corta y, en el
régimen postransitorio, nube, amplitudes, frecuencia y decisión de contacto con
tolerancias predeclaradas. La coincidencia de etiquetas controla la
clasificación, pero no estima por sí sola el orden de convergencia. Las variaciones de \(X_0\), \(h\), horizonte y memoria se evaluarán por separado.
En el caso con arctangente, la selección cruzada entre pasos descrita en
\cref{sec:refinamiento_c590} precede al sondeo de cuenca y constituye un
control distinto.
